\documentclass[11pt,a4paper]{article}

% ============== ENCODING AND FONTS ==============
\usepackage[utf8]{inputenc}
\usepackage[T1]{fontenc}
\usepackage{lmodern}
\usepackage{microtype}

% ============== PAGE LAYOUT ==============
\usepackage[margin=1in]{geometry}
\usepackage{setspace}
\onehalfspacing

% ============== MATHEMATICS ==============
\usepackage{amsmath,amssymb,amsfonts,amsthm,bm}
\usepackage{siunitx}
\usepackage{physics}
% Fix siunitx/physics clash: physics redefines \qty
\AtBeginDocument{\RenewCommandCopy\qty\SI}

% ============== GRAPHICS AND TABLES ==============
\usepackage{graphicx}
\usepackage{booktabs}
\usepackage{float}
\usepackage{caption}
\usepackage{subcaption}

% ============== LISTS AND ENUMERATIONS ==============
\usepackage{enumitem}

% ============== COLORS ==============
\usepackage{xcolor}

% ============== HYPERLINKS (load last) ==============
\usepackage[colorlinks=true,linkcolor=blue,citecolor=blue,urlcolor=blue]{hyperref}

% ============== CUSTOM MACROS ==============
\newcommand{\ee}[1]{\times 10^{#1}}
\newcommand{\kB}{\ensuremath{k_{\mathrm{B}}}}
\newcommand{\lPl}{\ensuremath{\ell_{\mathrm{Pl}}}}
\newcommand{\mPl}{\ensuremath{m_{\mathrm{Pl}}}}
\newcommand{\tPl}{\ensuremath{t_{\mathrm{Pl}}}}
\newcommand{\rhocrit}{\ensuremath{\rho_{\mathrm{crit}}}}
\newcommand{\rhostiff}{\ensuremath{\rho_{\mathrm{stiff}}}}
\newcommand{\rhoPl}{\ensuremath{\rho_{\mathrm{Pl}}}}
\newcommand{\Tzero}{\ensuremath{T_0}}
\newcommand{\Kbulk}{\ensuremath{K}}
\newcommand{\cs}{\ensuremath{c_s}}
\newcommand{\RH}{\ensuremath{R_{\mathrm{H}}}}
\newcommand{\umech}{\ensuremath{u_{\mathrm{mech}}}}
\newcommand{\wmech}{\ensuremath{w_{\mathrm{mech}}}}

% Paper X specific macros
\newcommand{\Teff}{\ensuremath{T_{\mathrm{eff}}}}
\newcommand{\Tthaw}{\ensuremath{T_{\mathrm{thaw}}}}
\newcommand{\TCMB}{\ensuremath{T_{\mathrm{CMB}}}}
\newcommand{\Rfc}{\ensuremath{R_{\mathrm{fc}}}}
\newcommand{\Mcross}{\ensuremath{M_{\mathrm{cross}}}}
\newcommand{\Zfc}{\ensuremath{Z_{\mathrm{fc}}}}
\newcommand{\Zpert}{\ensuremath{Z_{\mathrm{pert}}}}
\newcommand{\TH}{\ensuremath{T_{\mathrm{H}}}}
\newcommand{\rhofc}{\ensuremath{\rho_{\mathrm{fc}}}}

% TODO marker
\newcommand{\TODO}[1]{\textcolor{red}{\textbf{[TODO: #1]}}}
\newcommand{\NOTE}[1]{\textcolor{orange}{\textbf{[NOTE: #1]}}}

% Theorem environments
\newtheorem{theorem}{Theorem}[section]
\newtheorem{proposition}[theorem]{Proposition}
\newtheorem{corollary}[theorem]{Corollary}
\newtheorem{lemma}[theorem]{Lemma}
\theoremstyle{definition}
\newtheorem{definition}[theorem]{Definition}
\newtheorem{postulate}{Postulate}
\theoremstyle{remark}
\newtheorem{remark}[theorem]{Remark}
\newtheorem{example}[theorem]{Example}

% Proof QED symbol
\renewcommand{\qedsymbol}{$\blacksquare$}

% ============== TITLE ==============
\title{\textbf{Thermodynamics of Frozen Matter}\\[0.5em]
\large From Maximum Planck Density to Resting Tension: Absolute Zero of the Field, Thermal Thaw Pathway, and the Frozen Core Boundary}

\author{Robin Skavberg\\[0.3em]
\small Independent Researcher\\
\small ORCID: 0009-0003-9126-6196\\
\small \texttt{skavberg@gmail.com}}

\date{\today\\[0.5em]
\small Draft Version 01}

\begin{document}

\maketitle

\begin{abstract}
We develop the thermodynamic theory of frozen matter at maximum Planck field density $\rhostiff = c^2/(8\pi G) \approx 5.36 \times 10^{25}$ kg/m$^3$ within the Planck field condensate framework. At this density, matter reaches its ground state: all phonon modes cease, wave propagation becomes impossible, and the system exhibits zero thermodynamic excitations---a state we identify as the \emph{absolute zero of the Planck field}. This represents a fundamentally different thermodynamic limit than conventional absolute zero ($T = 0$ K at finite density), where quantum zero-point motion persists. Frozen matter at $\rhostiff$ corresponds to the resting tension $\Tzero = c^4/(8\pi G) \approx 4.83 \times 10^{42}$ Pa, the intrinsic elastic modulus of the field itself.

\vspace{0.5em}
\noindent
We derive the \emph{thermal thaw pathway}: the process by which frozen cores return to thermal equilibrium as their density decreases below $\rhostiff$. As the density relaxes, phonon modes progressively re-emerge, enabling thermodynamic excitations. We establish a density-temperature relation $\Teff(\rho)$ connecting the frozen state to the cosmic microwave background temperature $\TCMB = 2.725$ K at cosmological densities. The total energy budget for complete thaw equals the rest mass equivalent $Mc^2$, partitioned between thermal radiation and mechanical work during decompression.

\vspace{0.5em}
\noindent
The frozen core boundary exhibits a sharp thermal transition: the interior remains at $\Teff = 0$, while the exterior region where density falls below $\rhostiff$ begins to thermalize. We calculate the temperature profile, thermal gradient, heat flux, and luminosity of this boundary layer. For frozen cores with mass $M > \Mcross \approx 0.6\, M_{\mathrm{Jupiter}}$, this boundary radiation is redshifted by the gravitational potential, reproducing Hawking-like emission (Paper IX). For $M < \Mcross$, the boundary layer radiates directly into the surrounding medium without gravitational trapping.

\vspace{0.5em}
\noindent
Cosmological implications include primordial frozen cores thawing over Hubble timescales, contributing to diffuse photon backgrounds, and connecting to dark matter decay signatures. We provide testable predictions: characteristic emission spectra from frozen core boundaries, thermal evolution timescales, and constraints from observed X-ray, gamma-ray, and radio backgrounds. The thaw pathway provides the microscopic mechanism underlying the frozen core extrusion radiation described in Paper IX, unifying thermodynamics and gravity in the acoustic condensate framework.

\vspace{0.5em}
\noindent
\textbf{Keywords:} Planck field, frozen cores, thermodynamics, absolute zero, phonon dynamics, thermal thaw, dark matter stars, Hawking radiation, resting tension, acoustic black holes
\end{abstract}

\tableofcontents
\newpage

%==============================================================================
\section{Introduction}
\label{sec:introduction}
%==============================================================================

\subsection{Motivation: The Thermodynamic State at Maximum Density}
\label{sec:motivation}

The Planck field condensate framework (Paper I) establishes that matter at maximum compression reaches the stiffness-matching density $\rhostiff = c^2/(8\pi G) \approx 5.36 \times 10^{25}$ kg/m$^3$. At this density, the elastic modulus of compressed matter equals the resting tension of the Planck field $\Tzero = c^4/(8\pi G) \approx 4.83 \times 10^{42}$ Pa. This density represents a hard limit: no further compression is possible because the field itself cannot be made stiffer than its ground-state tension.

\vspace{0.5em}
\noindent
\textbf{Critical Distinction:} This paper addresses the thermodynamics of \emph{frozen cores}---compact objects of matter compressed to $\rhostiff$---not the cosmological Bose-Einstein condensate that constitutes the observable universe. The universe-as-BEC framework (superfluid dark matter, phonon dynamics, acoustic metric) provides the substrate in which frozen cores exist, but frozen cores themselves are localized, high-density structures distinct from the diffuse BEC background.

\vspace{0.5em}
\noindent
At $\rhostiff$, matter exhibits extraordinary properties:
\begin{enumerate}[leftmargin=*, itemsep=0.3em]
\item \textbf{Acoustic opacity:} All wave modes cease. Sound speed $\cs \to 0$ as the medium becomes incompressible.
\item \textbf{Zero excitations:} No phonon modes exist. The field is at its ground state with no available internal degrees of freedom.
\item \textbf{Thermodynamic ground state:} No thermal energy can be stored or transported. Entropy is minimized (possibly zero).
\item \textbf{Gravitational stability:} Frozen cores resist collapse because they have reached maximum density. No singularity forms (Paper IX).
\end{enumerate}

We propose that this state represents the \emph{absolute zero of the Planck field}---a thermodynamic limit fundamentally distinct from conventional absolute zero temperature.

\subsection{Absolute Zero of the Field vs. Conventional Absolute Zero}
\label{sec:two_zeros}

Conventional absolute zero ($T = 0$ K) corresponds to the quantum ground state of a system at finite density, where thermal excitations vanish but quantum zero-point motion persists. For example, liquid helium at $T \to 0$ K still exhibits zero-point fluctuations and superfluidity.

\vspace{0.5em}
\noindent
In contrast, the absolute zero of the Planck field occurs at maximum density $\rhostiff$, where:
\begin{itemize}[leftmargin=*, itemsep=0.3em]
\item \textbf{All wave propagation ceases:} $\cs = 0$ implies no phonon dispersion relation.
\item \textbf{No thermodynamic degrees of freedom:} Energy cannot be stored in internal modes.
\item \textbf{The field is maximally stiff:} Further compression is prohibited by the field's constitutive limit.
\end{itemize}

\TODO{Develop formal connection to third law of thermodynamics. Does $S \to 0$ at $\rhostiff$? Is this a violation or a refinement of Nernst's postulate?}

\NOTE{Compare with other ``zero'' concepts: Fermi energy at $T=0$, BEC condensation temperature, vacuum energy density. How does $\rhostiff$ relate to each?}

\subsection{The Thaw Problem: From Frozen Ground State to Thermal Equilibrium}
\label{sec:thaw_problem}

Given that frozen matter at $\rhostiff$ has zero effective temperature and zero phonon modes, how does it return to thermal equilibrium with its surroundings? This is the \emph{thermal thaw problem}.

\vspace{0.5em}
\noindent
Consider a frozen core of mass $M$ and radius $\Rfc = (6GM/c^2)^{1/3}$ embedded in the Planck field condensate at cosmological density $\rhocrit \approx 8.6 \times 10^{-27}$ kg/m$^3$. The surrounding medium is at temperature $\TCMB = 2.725$ K. As the frozen core slowly relaxes---either through quantum tunneling, boundary layer erosion, or gradual expansion---its density decreases below $\rhostiff$. At what point do phonon modes re-emerge? How does thermal energy redistribute? What is the end state?

\vspace{0.5em}
\noindent
We will show that the thaw pathway exhibits:
\begin{enumerate}[leftmargin=*, itemsep=0.3em]
\item \textbf{Progressive phonon activation:} As $\rho$ decreases below $\rhostiff$, low-frequency modes appear first, followed by higher-frequency modes as $\cs$ increases.
\item \textbf{Thermal energy release:} The rest mass energy $Mc^2$ is partitioned into thermal radiation (photons) and mechanical work (expansion against external pressure).
\item \textbf{Boundary layer dynamics:} The frozen core surface acts as a thermal barrier. Heat flux across this boundary determines the thaw rate.
\item \textbf{Final equilibrium:} Complete thaw occurs when the matter reaches cosmological density $\rhocrit$ and thermalizes to $\TCMB$.
\end{enumerate}

\TODO{Derive the phonon dispersion relation $\omega(k)$ as a function of density. At what $\rho/\rhostiff$ do modes become gapless?}

\subsection{Connection to Hawking Radiation (Paper IX)}
\label{sec:connection_hawking}

Paper IX establishes that frozen cores with $M > \Mcross \approx 0.6\, M_{\mathrm{Jupiter}}$ lie within an acoustic horizon, mimicking black hole phenomenology. The observed ``Hawking radiation'' is not quantum vacuum fluctuation, but \emph{extrusion radiation}---phononic excitations squeezed out from the frozen core boundary as matter is compressed toward $\rhostiff$.

\vspace{0.5em}
\noindent
The thermal thaw pathway provides the microscopic mechanism for this extrusion process:
\begin{itemize}[leftmargin=*, itemsep=0.3em]
\item Infalling matter approaches the frozen core boundary and experiences increasing density.
\item As $\rho \to \rhostiff$, phonon modes are progressively suppressed (the \emph{reverse thaw}).
\item The enormous acoustic impedance mismatch $\Zfc/\Zpert \sim 10^{52}$ reflects nearly all radiation, but the transmitted fraction produces observable emission.
\item This boundary radiation is redshifted by the gravitational potential, yielding a thermal spectrum with effective temperature $\TH$.
\end{itemize}

\NOTE{The thaw pathway is reversible: compression from $\rhocrit$ to $\rhostiff$ is the freezing process; relaxation from $\rhostiff$ to $\rhocrit$ is the thaw.}

\subsection{Paper Structure and Key Results}
\label{sec:structure}

This paper is organized as follows:

\begin{itemize}[leftmargin=*, itemsep=0.3em]
\item \textbf{Section~\ref{sec:ground_state}:} Formal characterization of the thermodynamic ground state at $\rhostiff$. Zero entropy, zero temperature, connection to third law.
\item \textbf{Section~\ref{sec:thaw_pathway}:} Derivation of the thermal thaw pathway: density-temperature relation $\Teff(\rho)$, phonon mode re-emergence, energy partition.
\item \textbf{Section~\ref{sec:reverse_calc}:} Reverse calculation from maximum density to CMB: total energy budget, $T(\rho)$ profile, comparison with Hawking temperature.
\item \textbf{Section~\ref{sec:boundary}:} Temperature profile at the frozen core boundary: thermal gradient, heat flux, luminosity, observational signatures.
\item \textbf{Section~\ref{sec:cosmology}:} Cosmological implications: primordial frozen cores, diffuse backgrounds, dark matter decay, observational constraints.
\item \textbf{Section~\ref{sec:connection_paper9}:} Connection to Hawking radiation (Paper IX): thaw as the microscopic mechanism, mass-dependent emission regimes.
\item \textbf{Section~\ref{sec:predictions}:} Testable predictions and falsification criteria.
\item \textbf{Section~\ref{sec:conclusion}:} Summary and outlook.
\end{itemize}

%==============================================================================
\section{Thermodynamic Ground State at Maximum Density}
\label{sec:ground_state}
%==============================================================================

\subsection{Density and Pressure at $\rhostiff$}
\label{sec:density_pressure}

From Paper I, the stiffness-matching density is:
\begin{equation}
\rhostiff = \frac{c^2}{8\pi G} \approx 5.36 \times 10^{25} \text{ kg/m}^3.
\end{equation}
This corresponds to the resting tension:
\begin{equation}
\Tzero = \rhostiff c^2 = \frac{c^4}{8\pi G} \approx 4.83 \times 10^{42} \text{ Pa}.
\end{equation}

At this density, the pressure-density relation reaches its maximum stiffness. The bulk modulus $\Kbulk = \rho \, \partial P/\partial \rho$ equals $\rhostiff c^2$, and further compression is prohibited.

\TODO{Derive the equation of state $P(\rho)$ explicitly. What functional form connects $\rho = \rhocrit$ to $\rho = \rhostiff$? Polytropic? Logarithmic? Empirical fit to numerical simulations?}

\subsection{Sound Speed and Acoustic Impedance}
\label{sec:sound_speed}

The sound speed is defined as:
\begin{equation}
\cs^2 = \frac{\partial P}{\partial \rho}.
\end{equation}
As $\rho \to \rhostiff$, the equation of state stiffens until $\cs \to 0$ at maximum density. This represents acoustic opacity: no wave can propagate through matter at $\rhostiff$.

\vspace{0.5em}
\noindent
The acoustic impedance is:
\begin{equation}
Z = \rho \cs.
\end{equation}
At the frozen core boundary:
\begin{equation}
\Zfc = \rhostiff \cdot 0 = \text{undefined (formally infinite in the limit)}.
\end{equation}

\TODO{Resolve the $0 \times \infty$ indeterminacy. What is the proper limiting form of $Z$ as $\rho \to \rhostiff$? Use $\cs \sim (\rhostiff - \rho)^{1/2}$ near the critical point?}

\subsection{Phonon Dispersion and Mode Suppression}
\label{sec:phonon_dispersion}

Phonons in a BEC condensate obey the Bogoliubov dispersion relation:
\begin{equation}
\omega(k) = c_s k \sqrt{1 + \frac{\xi^2 k^2}{2}},
\label{eq:bogoliubov}
\end{equation}
where $\xi = \hbar/(mc)$ is the healing length.

As $\cs \to 0$, all modes become gapped or vanish. At $\rhostiff$, no propagating modes exist: $\omega(k) = 0$ for all $k$. The system has no internal degrees of freedom capable of storing or transporting energy.

\TODO{Derive the mode density $g(\omega)$ as a function of $\rho$. How many modes are available at $\rho = 0.9\, \rhostiff$? At $\rho = 0.5\, \rhostiff$?}

\subsection{Entropy and the Third Law}
\label{sec:entropy}

The third law of thermodynamics (Nernst's postulate) states that the entropy of a perfect crystal approaches zero as $T \to 0$ K. Does this apply at $\rhostiff$?

\vspace{0.5em}
\noindent
If frozen matter at $\rhostiff$ has zero phonon modes, then:
\begin{equation}
S = \kB \ln \Omega = 0 \quad \text{if} \quad \Omega = 1.
\end{equation}
Here $\Omega$ is the number of accessible microstates. If the ground state is unique, $S = 0$ exactly.

\TODO{Investigate degeneracy. Could frozen cores have non-trivial topology or internal structure that increases $\Omega$? Spin configurations? Flavor degrees of freedom?}

\begin{proposition}[Zero Entropy at Maximum Density]
Matter at density $\rhostiff$ has zero entropy if and only if the ground state is non-degenerate.
\end{proposition}

\begin{proof}
\TODO{Complete proof. Requires explicit construction of the Hilbert space for matter at $\rhostiff$.}
\end{proof}

\subsection{Effective Temperature}
\label{sec:effective_temperature}

Define the effective temperature $\Teff$ such that:
\begin{equation}
\Teff = 0 \quad \text{at} \quad \rho = \rhostiff.
\end{equation}
This is not the kinetic temperature (which is undefined for a stationary condensate), but the thermodynamic temperature associated with excitations.

As $\rho$ decreases below $\rhostiff$, phonon modes re-emerge and $\Teff > 0$. The relationship $\Teff(\rho)$ is the central object of the thaw pathway analysis (Section~\ref{sec:thaw_pathway}).

\TODO{Propose functional form: $\Teff \propto (\rhostiff - \rho)^\alpha$ near $\rhostiff$? Or logarithmic? Exponential? Empirical fit?}

%==============================================================================
\section{The Thermal Thaw Pathway}
\label{sec:thaw_pathway}
%==============================================================================

\subsection{Starting Conditions: Frozen Matter at $\rhostiff$}
\label{sec:starting_conditions}

Consider a frozen core of mass $M$ at uniform density $\rhofc = \rhostiff$ and radius:
\begin{equation}
\Rfc = \left( \frac{6GM}{c^2} \right)^{1/3}.
\end{equation}
The core has:
\begin{itemize}[leftmargin=*, itemsep=0.3em]
\item Effective temperature $\Teff = 0$.
\item Zero entropy (assuming non-degenerate ground state).
\item Zero phonon modes.
\item Total rest energy $E_0 = Mc^2$.
\end{itemize}

The surrounding medium (Planck field BEC) is at density $\rhocrit$ and temperature $\TCMB = 2.725$ K.

\subsection{Density Relaxation and Phonon Re-Emergence}
\label{sec:phonon_reemergence}

As the frozen core expands or erodes, its density decreases below $\rhostiff$. At some critical density $\rho_{\mathrm{crit,phonon}}$, the first phonon modes become ungapped.

\vspace{0.5em}
\noindent
\textbf{Hypothesis:} Low-frequency acoustic modes re-emerge first. Higher-frequency modes require lower density (larger $\cs$) to propagate.

\TODO{Derive $\rho_{\mathrm{crit,phonon}}$ explicitly. Use Bogoliubov dispersion relation (Eq.~\ref{eq:bogoliubov}) and solve for the density at which $\omega(k_{\min}) > 0$.}

\subsection{Density-Temperature Relation}
\label{sec:density_temperature}

Propose the relation:
\begin{equation}
\Teff(\rho) = \TCMB \left( \frac{\rhostiff - \rho}{\rhostiff - \rhocrit} \right)^\alpha,
\label{eq:Teff_rho}
\end{equation}
where $\alpha$ is an exponent to be determined (possibly $\alpha = 1/2$ from dimensional analysis, or empirically fitted).

\vspace{0.5em}
\noindent
Boundary conditions:
\begin{align}
\Teff(\rhostiff) &= 0, \\
\Teff(\rhocrit) &= \TCMB.
\end{align}

\TODO{Derive $\alpha$ from first principles. Connect to phonon density of states, heat capacity, and Debye model.}

\subsection{Energy Partition: Radiation vs. Mechanical Work}
\label{sec:energy_partition}

The total rest energy $Mc^2$ is released as the frozen core thaws. This energy is partitioned into:
\begin{enumerate}[leftmargin=*, itemsep=0.3em]
\item \textbf{Thermal radiation:} Photons emitted from the boundary layer.
\item \textbf{Mechanical work:} PdV work done by the expanding core against the external pressure of the Planck field.
\end{enumerate}

Energy conservation:
\begin{equation}
Mc^2 = E_{\mathrm{rad}} + W_{\mathrm{mech}}.
\label{eq:energy_partition}
\end{equation}

\TODO{Derive the partition ratio. What fraction goes into radiation vs. mechanical work? Does this depend on $M$?}

\subsection{Thaw Timescale}
\label{sec:thaw_timescale}

The thaw rate depends on:
\begin{itemize}[leftmargin=*, itemsep=0.3em]
\item Thermal conductivity across the frozen core boundary.
\item Heat flux $q = -\kappa \nabla T$, where $\kappa$ is the thermal conductivity.
\item Boundary layer thickness $\delta \sim \Rfc (\Delta \rho/\rhostiff)$.
\end{itemize}

Estimate the total thaw time:
\begin{equation}
\tau_{\mathrm{thaw}} \sim \frac{Mc^2}{L_{\mathrm{boundary}}},
\label{eq:tau_thaw}
\end{equation}
where $L_{\mathrm{boundary}}$ is the luminosity of the boundary layer (Section~\ref{sec:boundary}).

\TODO{Compute $\tau_{\mathrm{thaw}}$ for representative masses: $M = 10^{30}$ kg (stellar mass), $M = 10^{20}$ kg (asteroid), $M = 10^{10}$ kg (mountain).}

%==============================================================================
\section{Reverse Calculation: From Maximum Density to CMB}
\label{sec:reverse_calc}
%==============================================================================

\subsection{The Thermodynamic Path}
\label{sec:thermo_path}

Consider the thermodynamic trajectory:
\begin{equation}
\text{Initial state:} \quad (\rho = \rhostiff, \, \Teff = 0) \quad \to \quad \text{Final state:} \quad (\rho = \rhocrit, \, T = \TCMB).
\end{equation}

The total entropy change is:
\begin{equation}
\Delta S = S_{\mathrm{final}} - S_{\mathrm{initial}} = S(\rhocrit, \TCMB) - 0.
\end{equation}

\TODO{Calculate $S(\rhocrit, \TCMB)$. Use ideal gas entropy formula? Photon gas entropy? BEC entropy formula?}

\subsection{Temperature Profile $T(\rho)$}
\label{sec:T_rho_profile}

From Eq.~\ref{eq:Teff_rho}, plot $\Teff(\rho)$ for $\rho$ ranging from $\rhostiff$ to $\rhocrit$. This gives the thermal evolution during thaw.

\TODO{Generate numerical plot. Include logarithmic density axis to span $10^{25}$ to $10^{-27}$ kg/m$^3$.}

\subsection{Energy Budget}
\label{sec:energy_budget}

The total energy released is exactly $Mc^2$. Compute the breakdown:
\begin{align}
E_{\mathrm{rad}} &= \int_{\rhostiff}^{\rhocrit} L_{\mathrm{boundary}}(\rho) \, dt, \\
W_{\mathrm{mech}} &= \int_{V_{\mathrm{initial}}}^{V_{\mathrm{final}}} P \, dV.
\end{align}

\TODO{Derive $L_{\mathrm{boundary}}(\rho)$ from heat flux calculation (Section~\ref{sec:boundary}). Integrate over the thaw trajectory.}

\subsection{Comparison with Hawking Temperature}
\label{sec:hawking_comparison}

For frozen cores with $M > \Mcross$, the boundary radiation is redshifted by the gravitational potential. The observed temperature is:
\begin{equation}
\TH = \frac{\hbar c^3}{8\pi G M \kB} \approx 6 \times 10^{-8} \, \text{K} \, \left( \frac{M_\odot}{M} \right).
\end{equation}

Compare $\TH$ with $\Teff(\rho)$ evaluated at the frozen core boundary. Are they consistent?

\TODO{Derive the redshift factor explicitly. Show that $\Teff$ at the boundary, when redshifted by the acoustic horizon geometry, yields $\TH$.}

%==============================================================================
\section{Temperature Profile at the Frozen Core Boundary}
\label{sec:boundary}
%==============================================================================

\subsection{Boundary Layer Structure}
\label{sec:boundary_structure}

The frozen core boundary is not a sharp discontinuity but a thin transition layer where:
\begin{itemize}[leftmargin=*, itemsep=0.3em]
\item Density decreases from $\rhostiff$ (interior) to $\rho < \rhostiff$ (exterior).
\item Effective temperature increases from $\Teff = 0$ (interior) to $\Teff > 0$ (exterior).
\item Phonon modes transition from fully suppressed to partially active.
\end{itemize}

The boundary layer thickness $\delta$ is estimated by:
\begin{equation}
\delta \sim \Rfc \left( \frac{\Delta \rho}{\rhostiff} \right),
\end{equation}
where $\Delta \rho \sim \rhostiff - \rho_{\mathrm{surface}}$.

\TODO{Estimate $\delta$ numerically. For $M = M_\odot$, what is $\delta$ in meters?}

\subsection{Thermal Gradient}
\label{sec:thermal_gradient}

The temperature gradient across the boundary is:
\begin{equation}
\frac{dT}{dr} \bigg|_{\mathrm{boundary}} = \frac{\Teff(r_{\mathrm{out}}) - 0}{\delta}.
\end{equation}

This gradient drives heat flux from the boundary into the surrounding medium.

\TODO{Compute $dT/dr$ for representative cases. How steep is this gradient?}

\subsection{Heat Flux and Luminosity}
\label{sec:heat_flux}

The heat flux is given by Fourier's law:
\begin{equation}
q = -\kappa \frac{dT}{dr},
\end{equation}
where $\kappa$ is the thermal conductivity of the Planck field near the boundary.

The total luminosity is:
\begin{equation}
L_{\mathrm{boundary}} = 4\pi \Rfc^2 \, q.
\label{eq:L_boundary}
\end{equation}

\TODO{Estimate $\kappa$. Use phonon thermal conductivity formula? Kinetic theory? Empirical fit from simulations?}

\TODO{Compute $L_{\mathrm{boundary}}$ for a range of masses. Compare with Hawking luminosity $L_{\mathrm{H}} = \hbar c^6/(15360\pi G^2 M^2)$.}

\subsection{Observational Signatures}
\label{sec:observational_signatures}

The boundary radiation spectrum is approximately thermal (blackbody) with temperature $\Teff(r_{\mathrm{boundary}})$. Observable features include:
\begin{enumerate}[leftmargin=*, itemsep=0.3em]
\item \textbf{Peak wavelength:} Wien's law gives $\lambda_{\max} = b/\Teff$.
\item \textbf{Total power:} Stefan-Boltzmann law gives $L \propto \Teff^4 \cdot A_{\mathrm{surface}}$.
\item \textbf{Spectral shape:} Planck function $B_\nu(T)$.
\end{enumerate}

\TODO{Generate synthetic spectra for frozen cores with $M = 10^{30}$ kg, $10^{20}$ kg, $10^{10}$ kg. What wavelengths dominate? Radio? Infrared? Optical? X-ray?}

%==============================================================================
\section{Cosmological Implications}
\label{sec:cosmology}
%==============================================================================

\subsection{Primordial Frozen Cores}
\label{sec:primordial_cores}

If frozen cores formed in the early universe (from density perturbations, phase transitions, or collapse of overdensities), they persist to the present day and contribute to dark matter.

\vspace{0.5em}
\noindent
The mass function $n(M)$ of primordial frozen cores depends on formation mechanisms and initial conditions. As these cores thaw over cosmological timescales $\tau_{\mathrm{thaw}} \sim t_{\mathrm{Hubble}}$, they release energy into the surrounding medium.

\TODO{Derive the contribution to the total energy density from thawing frozen cores. Compare with dark matter density $\Omega_{\mathrm{DM}} \approx 0.26$.}

\subsection{Contribution to Diffuse Backgrounds}
\label{sec:diffuse_backgrounds}

Thawing frozen cores emit photons across a range of wavelengths. The integrated emission contributes to:
\begin{itemize}[leftmargin=*, itemsep=0.3em]
\item Cosmic X-ray background (CXB).
\item Cosmic gamma-ray background (CGB).
\item Radio continuum backgrounds.
\item Infrared backgrounds.
\end{itemize}

\TODO{Compute the spectral intensity $I_\nu$ from a population of frozen cores. Integrate over mass function $n(M)$ and thaw timescale distribution.}

\subsection{Connection to Dark Matter Decay Signatures}
\label{sec:dm_decay}

If frozen cores are interpreted as dark matter candidates, their thaw process mimics dark matter decay or evaporation. Observational constraints from:
\begin{itemize}[leftmargin=*, itemsep=0.3em]
\item X-ray telescopes (Chandra, XMM-Newton, NuSTAR).
\item Gamma-ray observatories (Fermi-LAT, HESS, VERITAS).
\item Radio surveys (VLA, LOFAR, SKA precursors).
\end{itemize}

\TODO{Compare predicted emission rates with upper limits from searches for dark matter decay lines (e.g., 3.5 keV line).}

\subsection{Constraints from Observed Backgrounds}
\label{sec:constraints}

Existing measurements of diffuse backgrounds constrain the frozen core population. If too many cores thaw too rapidly, the predicted intensity exceeds observations.

\TODO{Derive upper limits on $n(M)$ from CXB and CGB measurements. What mass ranges are excluded?}

%==============================================================================
\section{Connection to Hawking Radiation (Paper IX)}
\label{sec:connection_paper9}
%==============================================================================

\subsection{Thaw as the Microscopic Mechanism}
\label{sec:thaw_mechanism}

Paper IX establishes that observed Hawking radiation from frozen cores is not quantum vacuum fluctuation, but \emph{extrusion radiation}---phononic excitations squeezed from the boundary as matter is compressed toward $\rhostiff$.

\vspace{0.5em}
\noindent
The thermal thaw pathway provides the reverse process: as frozen matter relaxes from $\rhostiff$, phonon modes re-emerge and radiate away. The compression (freezing) and relaxation (thaw) are thermodynamically reversible in principle, though dissipative in practice.

\TODO{Prove that the thaw pathway conserves energy and entropy along the trajectory. Show reversibility conditions.}

\subsection{Mass-Dependent Emission Regimes}
\label{sec:mass_regimes}

\subsubsection{Supercritical Frozen Cores ($M > \Mcross$)}
For $M > \Mcross \approx 0.6\, M_{\mathrm{Jupiter}}$, the frozen core lies within an acoustic horizon. Boundary radiation is redshifted by the gravitational potential. The observed temperature is:
\begin{equation}
\TH = \frac{\hbar c^3}{8\pi G M \kB}.
\end{equation}

This reproduces the Hawking temperature formula without invoking quantum vacuum fluctuations.

\subsubsection{Subcritical Frozen Cores ($M < \Mcross$)}
For $M < \Mcross$, no acoustic horizon exists. Boundary radiation escapes directly without gravitational redshift. The observed spectrum is determined by $\Teff$ at the boundary surface.

\TODO{Derive the transition criterion at $\Mcross$. What physical condition defines the acoustic horizon formation?}

\subsection{Information Preservation}
\label{sec:information}

The black hole information paradox arises because Hawking radiation is thermal and carries no information about the infalling matter. In the frozen core picture:
\begin{itemize}[leftmargin=*, itemsep=0.3em]
\item The interior is a finite-density condensate, not a singularity.
\item Information is stored in the condensate's quantum state.
\item Thaw radiation can, in principle, encode this information (though extracting it may be impractical).
\end{itemize}

\TODO{Formalize information content. Use von Neumann entropy? Entanglement entropy across the boundary?}

%==============================================================================
\section{Case Study: Thawing a Neutrino}
\label{sec:neutrino_thaw}
%==============================================================================

The preceding sections develop the thaw formalism in general terms. We now apply it to a specific particle---the neutrino---to obtain concrete numbers and physical insight. The neutrino is a natural choice: it is the lightest known massive particle, its mass hierarchy remains an open problem, and its weak interactions make it a minimal test case for the frozen core concept at particle physics scales.

\subsection{Neutrino Frozen Core Parameters}
\label{sec:neutrino_parameters}

From neutrino oscillation experiments, the mass eigenstates satisfy \cite{PDG2024}:
\begin{align}
\Delta m_{21}^2 &\approx 7.5 \times 10^{-5} \text{ eV}^2, \\
|\Delta m_{31}^2| &\approx 2.5 \times 10^{-3} \text{ eV}^2.
\end{align}
In the normal hierarchy ($m_1 < m_2 < m_3$), this gives approximate masses:
\begin{equation}
m_1 \approx 0, \quad m_2 \approx 0.009 \text{ eV}/c^2, \quad m_3 \approx 0.05 \text{ eV}/c^2.
\end{equation}
Cosmological constraints (Planck + BAO) bound the sum $\sum m_\nu < 0.12$ eV \cite{Planck2018cosmo}. We adopt $m_\nu = 0.05$ eV$/c^2 \approx 8.9 \times 10^{-38}$ kg for the heaviest eigenstate as our worked example.

\subsubsection{Frozen Core Radius}

At maximum density $\rhostiff$, the neutrino's mass is contained within a sphere of radius:
\begin{equation}
\Rfc = \left( \frac{3 m_\nu}{4\pi \rhostiff} \right)^{1/3} = \left( \frac{6 G m_\nu}{c^2} \right)^{1/3}.
\end{equation}
Substituting:
\begin{align}
\Rfc &= \left( \frac{6 \times 6.674 \times 10^{-11} \times 8.9 \times 10^{-38}}{(3 \times 10^8)^2} \right)^{1/3} \nonumber \\
     &= \left( 3.96 \times 10^{-64} \text{ m}^3 \right)^{1/3} \nonumber \\
     &\approx 7.3 \times 10^{-22} \text{ m}.
\end{align}

\begin{remark}[Scale Hierarchy]
The neutrino frozen core radius sits in a well-defined intermediate regime:
\begin{equation}
\lPl \approx 1.6 \times 10^{-35} \text{ m} \quad \ll \quad \Rfc \approx 7.3 \times 10^{-22} \text{ m} \quad \ll \quad r_p \approx 8.4 \times 10^{-16} \text{ m},
\end{equation}
where $\lPl$ is the Planck length and $r_p$ is the proton charge radius. The frozen core is $\sim 10^{13}$ times larger than the Planck length (so the semiclassical framework applies) but $\sim 10^{6}$ times smaller than a proton. This places neutrino frozen cores in the deep sub-nuclear regime, inaccessible to direct scattering experiments but within the domain where the Planck field framework operates.
\end{remark}

\subsubsection{Schwarzschild Radius Comparison}

The Schwarzschild radius of the neutrino is:
\begin{equation}
r_s = \frac{2 G m_\nu}{c^2} = \frac{2 \times 6.674 \times 10^{-11} \times 8.9 \times 10^{-38}}{9 \times 10^{16}} \approx 1.3 \times 10^{-64} \text{ m}.
\end{equation}
Since $r_s \ll \lPl \ll \Rfc$, the neutrino frozen core is vastly larger than its Schwarzschild radius. This confirms:
\begin{equation}
m_\nu \ll \Mcross \implies \text{no acoustic horizon forms.}
\end{equation}
The neutrino frozen core is \emph{subcritical}: it radiates directly during thaw without gravitational trapping or redshift.

\subsection{Thaw Energy Budget}
\label{sec:neutrino_energy}

The total energy released upon complete thaw equals the neutrino rest energy:
\begin{equation}
E_{\mathrm{thaw}} = m_\nu c^2 \approx 0.05 \text{ eV} \approx 8.0 \times 10^{-21} \text{ J}.
\end{equation}
This is the energy of a single infrared photon at wavelength $\lambda \approx hc/E \approx 25\,\mu$m.

\vspace{0.5em}
\noindent
From Eq.~\ref{eq:energy_partition}, this energy is partitioned:
\begin{equation}
m_\nu c^2 = E_{\mathrm{rad}} + W_{\mathrm{mech}}.
\end{equation}
For a subcritical frozen core with no acoustic horizon, the radiation is not gravitationally redshifted. The partition depends on the equation of state along the thaw trajectory.

\subsection{Compton Wavelength and the Frozen-to-Free Transition}
\label{sec:compton_transition}

The Compton wavelength of the neutrino is:
\begin{equation}
\lambda_C = \frac{\hbar}{m_\nu c} = \frac{1.055 \times 10^{-34}}{8.9 \times 10^{-38} \times 3 \times 10^8} \approx 4.0 \times 10^{-6} \text{ m} = 4\,\mu\text{m}.
\end{equation}

The ratio of Compton wavelength to frozen core radius quantifies the spatial expansion during thaw:
\begin{equation}
\frac{\lambda_C}{\Rfc} \approx \frac{4.0 \times 10^{-6}}{7.3 \times 10^{-22}} \approx 5.5 \times 10^{15}.
\end{equation}

\begin{remark}[Sixteen Orders of Magnitude]
\label{rem:sixteen_orders}
During the thaw, the neutrino's spatial extent expands by a factor of $\sim 10^{16}$---from a frozen core of radius $7.3 \times 10^{-22}$ m to a quantum wavepacket with Compton wavelength $4\,\mu$m. This enormous expansion is driven by the release of just $0.05$ eV of rest energy. The thaw is not a violent explosion but a gradual quantum delocalization as phonon modes re-emerge and the matter relaxes from $\rhostiff$ to its propagating state.
\end{remark}

\subsection{Thaw Dynamics at Particle Scales}
\label{sec:neutrino_dynamics}

As the neutrino frozen core begins to thaw, the density drops below $\rhostiff$ at the boundary. The thaw proceeds outward from the surface:

\begin{enumerate}[leftmargin=*, itemsep=0.3em]
\item \textbf{Phonon re-emergence:} The first modes to activate are long-wavelength (low-$k$) acoustic modes. For a frozen core of radius $\Rfc \sim 10^{-22}$ m, the longest mode has wavelength $\lambda_{\max} \sim 2\Rfc \sim 10^{-22}$ m, corresponding to frequency $\omega_{\max} \sim c_s / \Rfc$. As $c_s$ increases from zero, these modes become propagating.

\item \textbf{Decompression wave:} The boundary between frozen ($\rho = \rhostiff$) and thawing ($\rho < \rhostiff$) matter propagates inward at a speed determined by the local sound speed. For a neutrino-mass frozen core, the entire volume thaws once the decompression wave crosses the core diameter $\sim 2\Rfc$.

\item \textbf{Quantum delocalization:} As the density drops far below $\rhostiff$, the de Broglie wavelength of the matter exceeds the frozen core radius. The localized frozen core transitions into a delocalized quantum state---the free neutrino wavefunction. The crossover occurs when $\lambda_{\mathrm{dB}} \sim \Rfc$, i.e., when $\rho/\rhostiff \sim (\Rfc/\lambda_C)^3 \sim 10^{-48}$.
\end{enumerate}

\subsection{Thaw Timescale}
\label{sec:neutrino_timescale}

For a subcritical frozen core, the thaw timescale is bounded below by the light-crossing time:
\begin{equation}
\tau_{\min} = \frac{2\Rfc}{c} = \frac{2 \times 7.3 \times 10^{-22}}{3 \times 10^8} \approx 4.9 \times 10^{-30} \text{ s}.
\end{equation}
This is $\sim 10^{14}$ Planck times ($\tPl \approx 5.4 \times 10^{-44}$ s), confirming that the process is well above the Planck scale. The actual thaw time depends on the decompression speed, which is bounded by $c_s \leq c$.

\vspace{0.5em}
\noindent
For comparison, the natural timescale for neutrino oscillations is:
\begin{equation}
\tau_{\mathrm{osc}} = \frac{4\pi E}{\Delta m^2 c^4} \sim 10^{-11} \text{ s} \quad \text{(for $E \sim 1$ MeV)}.
\end{equation}
The thaw timescale $\tau_{\min} \sim 10^{-30}$ s is $\sim 10^{19}$ times shorter than the oscillation timescale. This suggests that frozen core formation and dissolution occur instantaneously relative to flavor dynamics.

\subsection{Mass Hierarchy as Frozen Core Hierarchy}
\label{sec:mass_hierarchy}

The three neutrino mass eigenstates correspond to three distinct frozen core configurations:
\begin{center}
\begin{tabular}{llll}
\toprule
\textbf{Eigenstate} & \textbf{Mass (eV/$c^2$)} & \textbf{$\Rfc$ (m)} & \textbf{$\lambda_C$ (m)} \\
\midrule
$\nu_1$ & $\sim 0.001$ & $\sim 2.0 \times 10^{-22}$ & $\sim 2.0 \times 10^{-4}$ \\
$\nu_2$ & $\sim 0.009$ & $\sim 4.1 \times 10^{-22}$ & $\sim 2.2 \times 10^{-5}$ \\
$\nu_3$ & $\sim 0.050$ & $\sim 7.3 \times 10^{-22}$ & $\sim 4.0 \times 10^{-6}$ \\
\bottomrule
\end{tabular}
\end{center}

The frozen core radius scales as $\Rfc \propto m_\nu^{1/3}$, so the mass hierarchy maps to a weak hierarchy in frozen core size. The heaviest eigenstate ($\nu_3$) has a frozen core $\sim 3.7\times$ larger than the lightest ($\nu_1$), since $(m_3/m_1)^{1/3} = (50)^{1/3} \approx 3.7$.

\begin{remark}[Flavor Mixing and Frozen Core Superposition]
In the standard model, flavor eigenstates ($\nu_e, \nu_\mu, \nu_\tau$) are superpositions of mass eigenstates ($\nu_1, \nu_2, \nu_3$). In the frozen core picture, a flavor eigenstate corresponds to a \emph{quantum superposition of frozen core configurations} with different radii and densities. Neutrino oscillations then correspond to the time evolution of this superposition, with the phase difference accumulating because the different frozen core configurations have different rest energies. This interpretation connects directly to Paper~XI (quantum superposition at maximum density).
\end{remark}

\subsection{Particle Mass as Thaw Energy}
\label{sec:mass_as_thaw}

The thaw framework suggests a physical interpretation of particle rest mass:

\begin{proposition}[Rest Mass as Thaw Energy]
\label{prop:mass_thaw}
The rest mass $m$ of any massive particle equals the total energy required to completely thaw a frozen core of that particle's mass from $\rhostiff$ to its propagating state:
\begin{equation}
mc^2 = \int_{\rhostiff}^{\rho_{\mathrm{free}}} \left[ \frac{\partial E_{\mathrm{rad}}}{\partial \rho} + P \frac{\partial V}{\partial \rho} \right] d\rho.
\end{equation}
\end{proposition}

This is a tautology by energy conservation (Eq.~\ref{eq:energy_partition}), but it suggests a deeper interpretation: \emph{rest mass is the energy cost of maintaining a localized field displacement against the Planck field's resting tension}. A neutrino's $0.05$ eV mass reflects the minimal energy needed to sustain its frozen core configuration.

\subsubsection{Standard Model Frozen Core Hierarchy}

The frozen core radius $\Rfc = (3m/(4\pi\rhostiff))^{1/3}$ scales as $m^{1/3}$, compressing the enormous mass hierarchy of the Standard Model into a more modest range of frozen core sizes. Table~\ref{tab:frozen_cores} lists all fundamental fermions, the Higgs and $W$ bosons, and the proton (as a composite reference) ordered by increasing mass.

\begin{table}[H]
\centering
\caption{Frozen core parameters for Standard Model particles and proto-particles. Radii computed from $\Rfc = (3m/(4\pi\rhostiff))^{1/3}$ with $\rhostiff = 5.36 \times 10^{25}$ kg/m$^3$. Quark masses are current (running) masses at $\mu = 2$ GeV \cite{PDG2024}. The Planck length is $\lPl = 1.616 \times 10^{-35}$ m.}
\label{tab:frozen_cores}
\vspace{0.5em}
\begin{tabular}{lllll}
\toprule
\textbf{Particle} & \textbf{Mass (eV/$c^2$)} & \textbf{$\Rfc$ (m)} & \textbf{$\Rfc/\lPl$} & \textbf{Notes} \\
\midrule
\multicolumn{5}{l}{\textit{Leptons}} \\
$\nu_3$ (neutrino) & $0.05$ & $7.3 \times 10^{-22}$ & $4.5 \times 10^{13}$ & Lightest massive particle \\
$e$ (electron) & $5.11 \times 10^5$ & $1.6 \times 10^{-19}$ & $1.0 \times 10^{16}$ & Lightest charged lepton \\
$\mu$ (muon) & $1.06 \times 10^8$ & $9.3 \times 10^{-19}$ & $5.8 \times 10^{16}$ & \\
$\tau$ (tau) & $1.78 \times 10^9$ & $2.4 \times 10^{-18}$ & $1.5 \times 10^{17}$ & \\
\midrule
\multicolumn{5}{l}{\textit{Quarks (current masses at $\mu = 2$ GeV)}} \\
$u$ (up) & $2.16 \times 10^6$ & $2.6 \times 10^{-19}$ & $1.6 \times 10^{16}$ & Confined; current mass \\
$d$ (down) & $4.67 \times 10^6$ & $3.3 \times 10^{-19}$ & $2.0 \times 10^{16}$ & Confined; current mass \\
$s$ (strange) & $9.34 \times 10^7$ & $9.0 \times 10^{-19}$ & $5.6 \times 10^{16}$ & \\
$c$ (charm) & $1.27 \times 10^9$ & $2.2 \times 10^{-18}$ & $1.4 \times 10^{17}$ & \\
$b$ (bottom) & $4.18 \times 10^9$ & $3.2 \times 10^{-18}$ & $2.0 \times 10^{17}$ & \\
$t$ (top) & $1.73 \times 10^{11}$ & $1.1 \times 10^{-17}$ & $6.8 \times 10^{17}$ & Heaviest SM fermion \\
\midrule
\multicolumn{5}{l}{\textit{Bosons}} \\
gluon $g$ & $0$ & $0$ & --- & Massless; see \S\ref{sec:gluon_thaw} \\
$W^{\pm}$ & $8.04 \times 10^{10}$ & $8.6 \times 10^{-18}$ & $5.3 \times 10^{17}$ & \\
$Z^0$ & $9.12 \times 10^{10}$ & $8.9 \times 10^{-18}$ & $5.5 \times 10^{17}$ & \\
$H^0$ (Higgs) & $1.25 \times 10^{11}$ & $1.0 \times 10^{-17}$ & $6.2 \times 10^{17}$ & \\
\midrule
\multicolumn{5}{l}{\textit{Composite (for reference)}} \\
$p$ (proton) & $9.38 \times 10^8$ & $2.0 \times 10^{-18}$ & $1.2 \times 10^{17}$ & 99\% gluon field energy \\
\midrule
\multicolumn{5}{l}{\textit{Proto-particles (hypothetical)}} \\
GUT-scale & $\sim 10^{25}$ & $4.3 \times 10^{-13}$ & $2.7 \times 10^{22}$ & Grand unification scale \\
Planck mass & $1.22 \times 10^{28}$ & $4.6 \times 10^{-12}$ & $2.8 \times 10^{23}$ & Maximum elementary $\Rfc$ \\
\bottomrule
\end{tabular}
\end{table}

All Standard Model frozen core radii satisfy $\Rfc \gg \lPl$ by at least 13 orders of magnitude, confirming that the semiclassical framework applies across the entire particle mass spectrum. The frozen core radius ranges from $7.3 \times 10^{-22}$ m (neutrino) to $1.1 \times 10^{-17}$ m (top quark)---a span of only $\sim 10^{4.2}$, compared to $\sim 10^{12.5}$ for the corresponding mass range.

\subsubsection{The Quark Sector: Current Mass vs. Constituent Mass}
\label{sec:quark_thaw}

The quark entries in Table~\ref{tab:frozen_cores} use the \emph{current masses} (running masses at $\mu = 2$ GeV), which reflect the bare Lagrangian parameters of QCD. These are distinct from \emph{constituent masses} ($m_u^{\mathrm{const}} \approx 336$ MeV, $m_d^{\mathrm{const}} \approx 340$ MeV), which include the dynamically generated mass from chiral symmetry breaking and gluon dressing.

\vspace{0.5em}
\noindent
In the frozen core picture, this distinction is significant:
\begin{itemize}[leftmargin=*, itemsep=0.3em]
\item \textbf{Current mass frozen core:} The bare quark, stripped of its gluon cloud, displaced the Planck field by only $2.16$ MeV (for the up quark). This is the intrinsic field displacement of the quark itself.
\item \textbf{Constituent mass frozen core:} The dressed quark, including its gluon field energy, displaced the field by $\sim 336$ MeV. The additional $\sim 334$ MeV comes from the gluon field configuration, not from quark rest mass.
\end{itemize}

Since quarks are confined and never observed in isolation, the physically relevant frozen core for a quark is embedded within a hadron. The thaw of a quark frozen core would therefore involve the simultaneous reconstitution of the confining gluon field---it cannot occur in isolation.

\subsubsection{Gluons: Massless Particles and the Binding Energy Problem}
\label{sec:gluon_thaw}

Gluons are massless ($m_g = 0$) and therefore have $\Rfc = 0$: they do not form frozen cores. However, gluon field energy dominates the mass budget of all hadrons. The proton's mass of $938$ MeV decomposes as:
\begin{equation}
m_p c^2 \approx \underbrace{9.4 \text{ MeV}}_{\text{quark masses}} + \underbrace{929 \text{ MeV}}_{\text{gluon field energy}}.
\end{equation}

This creates a striking asymmetry in the thaw framework:
\begin{itemize}[leftmargin=*, itemsep=0.3em]
\item The proton's frozen core has $\Rfc \approx 2.0 \times 10^{-18}$ m, computed from its \emph{total} mass $938$ MeV.
\item Only $\sim 1\%$ of the thaw energy goes into reconstituting quark rest masses.
\item The remaining $\sim 99\%$ reconstitutes the gluon field configuration---the chromodynamic binding that makes a proton a proton.
\end{itemize}

\begin{remark}[Mass Without Rest Mass]
\label{rem:gluon_mass}
The gluon contribution to hadron mass demonstrates that a frozen core's thaw energy need not originate from the rest masses of its constituents. A proton-mass frozen core at $\rhostiff$ stores $938$ MeV of field displacement energy. Upon thaw, most of this energy goes into creating a gluon field configuration (which itself has zero rest mass constituents) rather than into massive particles. This suggests that \emph{mass} in the Planck field framework is fundamentally about the magnitude of field displacement, not about intrinsic particle properties.
\end{remark}

\subsubsection{Proto-Particles: Protofluid Excitations at Extreme Scales}
\label{sec:proto_particles}

The Planck field framework postulates a universal protofluid into which the observable BEC condensate expands. Excitations of this protofluid at energy scales far above the Standard Model constitute \emph{proto-particles}---the elementary degrees of freedom of the pre-condensation medium.

\vspace{0.5em}
\noindent
Two characteristic energy scales define the proto-particle regime:

\textbf{Grand Unified Theory (GUT) scale} ($E \sim 10^{16}$ GeV $= 10^{25}$ eV):
\begin{itemize}[leftmargin=*, itemsep=0.2em]
\item Frozen core radius $\Rfc \sim 4.3 \times 10^{-13}$ m $\approx 0.43$ pm---comparable to atomic radii.
\item $\Rfc/\lPl \sim 10^{22}$: deeply semiclassical.
\item At GUT scales, the electromagnetic, weak, and strong forces are expected to unify. A GUT-scale frozen core would contain the unified force structure in a single compressed field displacement.
\item Thaw energy $\sim 10^{16}$ GeV per particle. If primordial GUT-scale frozen cores persist, their thaw would produce ultra-high-energy particles far beyond any accelerator reach.
\end{itemize}

\textbf{Planck scale} ($E \sim 1.22 \times 10^{19}$ GeV $= 1.22 \times 10^{28}$ eV):
\begin{itemize}[leftmargin=*, itemsep=0.2em]
\item Frozen core radius $\Rfc \sim 4.6 \times 10^{-12}$ m $\approx 4.6$ pm.
\item $\Rfc/\lPl \sim 10^{23}$: still within the semiclassical regime, which is notable---even at the Planck mass, the frozen core is vastly larger than the Planck length.
\item A Planck-mass frozen core represents the largest possible elementary field displacement. Beyond this scale, the frozen core becomes macroscopic and constitutes a dark compact object rather than a particle.
\item The Planck mass marks the boundary between particle physics and astrophysics in the frozen core hierarchy.
\end{itemize}

\begin{remark}[Frozen Core Radius at the Planck Mass]
The Planck-mass frozen core has radius $\Rfc(\mPl) = (3\mPl/(4\pi\rhostiff))^{1/3} \approx 4.6 \times 10^{-12}$ m, which is $\sim 10^{23}\,\lPl$. This large ratio shows that even at the highest conceivable particle energy, the frozen core formalism remains far from the quantum gravity regime where $\Rfc \sim \lPl$. The semiclassical approximation breaks down only for hypothetical objects with mass $m \sim \lPl^3 \rhostiff \sim 10^{-80}$ kg---a mass so small it is unphysical. This confirms that the frozen core framework is self-consistent across the entire particle spectrum.
\end{remark}

\subsubsection{The Mass Spectrum as a Frozen Core Spectrum}

The full particle table reveals a pattern: every massive particle corresponds to a specific frozen core radius, and the entire Standard Model mass spectrum maps onto a frozen core size spectrum spanning $7 \times 10^{-22}$ m to $1.1 \times 10^{-17}$ m. Extending to proto-particles pushes this to $\sim 5 \times 10^{-12}$ m.

\vspace{0.5em}
\noindent
The $\Rfc \propto m^{1/3}$ scaling compresses the mass hierarchy:
\begin{itemize}[leftmargin=*, itemsep=0.3em]
\item The top quark is $\sim 3 \times 10^{12}$ times heavier than the neutrino, but its frozen core is only $\sim 1.5 \times 10^4$ times larger.
\item The proton ($938$ MeV) and charm quark ($1.27$ GeV) have nearly identical frozen core radii ($2.0$ vs.\ $2.2 \times 10^{-18}$ m), despite the charm quark being 35\% heavier, because $\Rfc$ depends on $m^{1/3}$.
\end{itemize}

\TODO{Investigate whether the frozen core radius hierarchy reproduces any known mass relations (e.g., Koide formula for charged lepton masses, or the electron-proton mass ratio $m_p/m_e \approx 1836$). Explore whether the quantization of $\Rfc$ could explain the discrete particle mass spectrum.}

\subsection{Implications for Neutrino Cosmology}
\label{sec:neutrino_cosmology}

The neutrino thaw picture has implications for cosmological neutrino physics:

\begin{enumerate}[leftmargin=*, itemsep=0.3em]
\item \textbf{Cosmic neutrino background (C$\nu$B):} Relic neutrinos from the early universe ($T_\nu \approx 1.95$ K today) have momenta $p \sim 3\kB T_\nu/c \sim 5 \times 10^{-4}$ eV$/c$, comparable to their rest mass. In the frozen core picture, these neutrinos are in a partially thawed state---neither fully frozen nor fully delocalized. The transition between relativistic ($p \gg mc$) and non-relativistic ($p \ll mc$) regimes at $z \sim m_\nu c^2/(3\kB T_{\nu,0}) \sim 100$ corresponds to the onset of ``re-freezing'' in the Planck field language.

\item \textbf{Neutrino free-streaming scale:} Neutrinos suppress structure formation below their free-streaming scale $k_{\mathrm{fs}} \sim 0.018\,\Omega_m^{1/2}\,(m_\nu/1\,\text{eV})\,h/\text{Mpc}$ \cite{Lesgourgues2006}. In the frozen core picture, this scale reflects the spatial extent of the neutrino's thawed wavefunction---structure cannot form on scales smaller than the neutrino's quantum delocalization length.

\item \textbf{Neutrino mass from thaw energy:} If the thaw energy can be calculated from first principles (i.e., from the Planck field equation of state), this would predict the neutrino mass. This remains an open problem requiring the explicit form of $P(\rho)$ near $\rhostiff$.
\end{enumerate}

%==============================================================================
\section{Predictions and Falsification Criteria}
\label{sec:predictions}
%==============================================================================

The thermal thaw framework makes testable predictions:

\begin{enumerate}[leftmargin=*, itemsep=0.3em]
\item \textbf{Boundary Emission Spectrum:} Frozen cores emit approximately thermal radiation with peak wavelength determined by $\Teff$ at the boundary. For stellar-mass cores, expect X-ray emission; for planetary-mass cores, infrared or radio emission.

\item \textbf{Thaw Timescales:} The luminosity $L_{\mathrm{boundary}}$ determines $\tau_{\mathrm{thaw}} \sim Mc^2/L$. For primordial cores with $M \sim 10^{20}$ kg, predict $\tau_{\mathrm{thaw}} \sim 10^{10}$ yr (comparable to Hubble time).

\item \textbf{Contribution to Diffuse Backgrounds:} If a significant population of frozen cores exists, their collective emission contributes to CXB, CGB, or radio backgrounds. Upper limits from observations constrain $n(M)$.

\item \textbf{Spectral Features:} Unlike purely thermal Hawking radiation, boundary emission may exhibit line features or deviations from perfect blackbody if the boundary layer has non-trivial composition (e.g., mixed phases, density gradients).

\item \textbf{Temporal Evolution:} Thawing frozen cores should exhibit slow brightening over cosmological timescales as $L_{\mathrm{boundary}}$ increases with decreasing $\rho$.
\end{enumerate}

\vspace{0.5em}
\noindent
\textbf{Falsification Criteria:}
\begin{itemize}[leftmargin=*, itemsep=0.3em]
\item If observed diffuse backgrounds significantly exceed predictions, the frozen core population must be smaller than assumed.
\item If no isolated compact objects exhibit the predicted boundary emission spectrum, frozen cores may not exist as stable configurations.
\item If Hawking radiation from astrophysical black holes shows quantum information loss beyond the thermal limit, the frozen core picture is incomplete.
\end{itemize}

\TODO{Develop detailed observational strategies. Which telescopes? Which wavebands? What integration times?}

%==============================================================================
\section{Conclusion}
\label{sec:conclusion}
%==============================================================================

We have developed the thermodynamic theory of frozen matter at maximum Planck field density $\rhostiff = c^2/(8\pi G)$. At this density, matter reaches a ground state with zero phonon modes, zero effective temperature, and zero entropy---the \emph{absolute zero of the Planck field}. This is distinct from conventional absolute zero, which occurs at $T = 0$ K at finite density.

\vspace{0.5em}
\noindent
The \emph{thermal thaw pathway} describes how frozen cores return to thermal equilibrium as their density decreases below $\rhostiff$. Phonon modes progressively re-emerge, enabling thermodynamic excitations. The density-temperature relation $\Teff(\rho)$ connects the frozen state to the cosmic microwave background temperature $\TCMB = 2.725$ K at cosmological densities. The total energy budget equals the rest mass equivalent $Mc^2$, partitioned between thermal radiation and mechanical work.

\vspace{0.5em}
\noindent
The frozen core boundary exhibits a sharp thermal transition, with a temperature gradient driving heat flux and luminosity. For supercritical masses $M > \Mcross$, this boundary radiation is redshifted by the gravitational potential, reproducing Hawking-like emission (Paper IX). For subcritical masses, the radiation escapes directly without gravitational trapping.

\vspace{0.5em}
\noindent
Cosmological implications include primordial frozen cores thawing over Hubble timescales, contributing to diffuse photon backgrounds, and connecting to dark matter decay signatures. We have provided testable predictions: characteristic emission spectra, thaw timescales, and constraints from observed backgrounds.

\vspace{0.5em}
\noindent
The thaw pathway unifies thermodynamics and gravity in the Planck field condensate framework. It provides the microscopic mechanism underlying frozen core extrusion radiation, resolves the information paradox by eliminating singularities, and offers a concrete alternative to black hole thermodynamics based on acoustic geometry rather than quantum vacuum fluctuations.

\vspace{0.5em}
\noindent
\textbf{Future Work:}
\begin{itemize}[leftmargin=*, itemsep=0.3em]
\item Numerical simulations of the thaw process using Gross-Pitaevskii hydrodynamics.
\item Detailed spectral modeling of boundary emission for comparison with observations.
\item Connection to gravitational wave signatures from merging frozen cores (Paper VI).
\item Quantum superposition of frozen core states near $\rhostiff$ (Paper XI).
\end{itemize}

%==============================================================================
\section*{Acknowledgments}
%==============================================================================

This work is part of the Planck Field BEC Cosmology research program. The author acknowledges helpful discussions on thermodynamics of condensed matter systems and acoustic black hole phenomenology.

%==============================================================================
\begin{thebibliography}{99}
%==============================================================================

\bibitem{Paper1}
Skavberg, R. (2025).
\textit{The Observable Universe as a Condensate.}
Draft Version 03.

\bibitem{Paper5}
Skavberg, R. (2025).
\textit{Dark Energy as Protofluid Negative Pressure.}
In preparation.

\bibitem{Paper9}
Skavberg, R. (2025).
\textit{Hawking Radiation as Frozen Core Extrusion.}
Draft Version 01.

\bibitem{PDG2024}
R.~L.~Workman et al. (Particle Data Group) (2024).
\textit{Review of Particle Physics.}
Progress of Theoretical and Experimental Physics, 2024, 083C01.

\bibitem{Planck2018cosmo}
N.~Aghanim et al. (Planck Collaboration) (2020).
\textit{Planck 2018 Results. VI. Cosmological Parameters.}
Astronomy \& Astrophysics, 641, A6.

\bibitem{Lesgourgues2006}
J.~Lesgourgues and S.~Pastor (2006).
\textit{Massive Neutrinos and Cosmology.}
Physics Reports, 429(6), 307--379.

\bibitem{Jacobson1995}
Jacobson, T. (1995).
\textit{Thermodynamics of Spacetime: The Einstein Equation of State.}
Physical Review Letters, 75(7), 1260.

\bibitem{Padmanabhan2010}
Padmanabhan, T. (2010).
\textit{Thermodynamical Aspects of Gravity: New insights.}
Reports on Progress in Physics, 73(4), 046901.

\bibitem{Unruh1981}
Unruh, W. G. (1981).
\textit{Experimental Black-Hole Evaporation?}
Physical Review Letters, 46(21), 1351.

\bibitem{Barcelo2005}
Barceló, C., Liberati, S., \& Visser, M. (2005).
\textit{Analogue Gravity.}
Living Reviews in Relativity, 8(1), 12.

\bibitem{Berezhiani2015}
Berezhiani, L., \& Khoury, J. (2015).
\textit{Theory of Dark Matter Superfluidity.}
Physical Review D, 92(10), 103510.

\bibitem{Suarez2014}
Suárez, A., \& Chavanis, P.-H. (2014).
\textit{Cosmological Evolution of a Complex Scalar Field with Repulsive or Attractive Self-Interaction.}
Physical Review D, 92(2), 023510.

\bibitem{Sharma2020}
Sharma, R., et al. (2020).
\textit{Ultralight Dark Matter in Bose-Einstein Condensates.}
Physics Reports, 882, 1-76.

\bibitem{Fermi1936}
Fermi, E. (1936).
\textit{Thermodynamics.}
Dover Publications.

\bibitem{Landau1980}
Landau, L. D., \& Lifshitz, E. M. (1980).
\textit{Statistical Physics, Part 1.}
Pergamon Press.

\bibitem{Pethick2008}
Pethick, C. J., \& Smith, H. (2008).
\textit{Bose-Einstein Condensation in Dilute Gases.}
Cambridge University Press.

\bibitem{Pitaevskii2016}
Pitaevskii, L., \& Stringari, S. (2016).
\textit{Bose-Einstein Condensation and Superfluidity.}
Oxford University Press.

\end{thebibliography}

\end{document}
